\documentclass[12pt]{article}
%---DOCUMENT MARGINS---
\usepackage{geometry} % Required for adjusting page dimensions and margins
\geometry{
	paper=a4paper, % Paper size, change to letterpaper for US letter size
	top=4cm, % Top margin
	bottom=2cm, % Bottom margin
	left=2cm, % Left margin
	right=2cm, % Right margin
	headheight=3cm, % Header height
	%footskip=1.5cm, % Space from the bottom margin to the baseline of the footer
	headsep=0.5cm, % Space from the top margin to the baseline of the header
	%showframe, % Uncomment to show how the type block is set on the page
}
\usepackage{cmap}

\usepackage[T1,T2A]{fontenc}
\usepackage{fontspec}
\setmainfont[
	BoldFont={* Bold},
	ItalicFont={* Italic},
	BoldItalicFont={* BoldItalic}
]{DejaVu Serif Condensed}
\setsansfont{DejaVu Sans}
\setmonofont{Dejavu Sans Mono}
\usepackage[math-style=TeX]{unicode-math}
\setmathfont{Latin Modern Math}

\usepackage[nobottomtitles*]{titlesec}
\titleformat
{\section} % command
{\fontsize{14}{14}\sffamily\bfseries} % format
{} % label
{0pt} % sep
{} % before-code
\titlespacing{\section}{0pt}{0.75em}{0.25em}
\newcommand{\tl}{}
\newcommand{\ml}{}
\newcommand{\problem}[1]{%
	\section[#1]{#1 \fontsize{14}{14}\selectfont{\hfill{\normalfont{
				\emoji{hourglass-not-done}\tl \space\space \emoji{floppy-disk} \ml}}}}
}
\AddToHook{cmd/section/before}{\clearpage}

\usepackage[dvipsnames]{xcolor}
\titleformat
{\subsection} % command
{\fontsize{14}{14}\sffamily\bfseries} % format
{\includesvg[width=0.035\textwidth, inkscapelatex=false]{lion}\space} % label
{0pt} % sep
{} % before-code
\titlespacing{\subsection}{0pt}{0.75em}{0.25em}

\setlength{\parskip}{0.5em}
\setlength{\parindent}{0pt}
\renewcommand{\baselinestretch}{1.2}
\sloppy

\usepackage{listings}
\definecolor{lstbg}{RGB}{250,250,252}
\definecolor{lstframe}{RGB}{220,220,225}
\lstdefinestyle{mystyle}{
	backgroundcolor=\color{lstbg},
	frame=single,
	rulecolor=\color{lstframe},
	basicstyle=\ttfamily\small,
	keywordstyle=\bfseries,
	breaklines=true,
	aboveskip=0.5\baselineskip,
	belowskip=0\baselineskip  
}
\lstset{style=mystyle, language=C++}

\usepackage{fancyhdr}
\pagestyle{fancy}
\setlength{\headwidth}{\textwidth}
\newcommand{\round}{}
\newcommand{\problemName}{}
\newcommand{\lang}{English}
\fancyhead[L]{
	\begin{minipage}{0.4\textwidth}
		\bf{
			EJOI 2025 \round\\
			\problemName\\
			\emoji{globe-with-meridians} \lang
		}
	\end{minipage}
	\vspace{0.5cm}
}
\fancyhead[R]{%
	\begin{minipage}{0.6\textwidth}
		\includesvg[width=\textwidth, inkscapelatex=false]{logo}
	\end{minipage}
	\vspace{0.5cm}
}
\usepackage{lastpage}
\fancyfoot[C]{\problemName\space(\lang) \thepage\ / \pageref*{LastPage}}

\raggedbottom

\usepackage{amsmath}
\usepackage{stmaryrd}

\usepackage{graphicx}
\graphicspath{{./}}
\usepackage[export]{adjustbox}
\usepackage{wrapfig}
\makeatletter
\patchcmd\WF@putfigmaybe{\lower\intextsep}{}{}{\fail}
\AddToHook{env/wrapfigure/begin}{\setlength{\intextsep}{0pt}}
\makeatother
\usepackage[inkscapearea=page, inkscapepath=./svg-inkscape]{svg}
\svgpath{{./}}

\usepackage{makecell}
\usepackage{tabularray}
\AtBeginEnvironment{table}{\vspace{-0.2cm}}
\AtEndEnvironment{table}{\vspace{-0.2cm}}
\usepackage{float}
\usepackage{placeins}
\usepackage{caption}
\captionsetup[table]{
	skip=0.25em,
	singlelinecheck=false, justification=justified
}

\usepackage{enumitem}
\setlist{itemsep=-0.4em, leftmargin=22pt, topsep=-\parskip}
\AfterEndEnvironment{itemize}{\vspace{\parskip}}
\newcommand{\tabitem}{\indent~~\llap{\textbullet}~~}

\usepackage{hyperref}
\hypersetup{
	colorlinks=true,
	citecolor=blue,
	linkcolor=blue,
	urlcolor=cyan,
}

\usepackage{emoji}

\usepackage[english]{babel}

\begin{document}

\renewcommand{\round}{Ден 2}
\renewcommand{\problemName}{Задача Навигација}
\renewcommand{\lang}{Македонски}

\renewcommand{\tl}{$3$ сек.}
\renewcommand{\ml}{$512$ MB}
\problem{\problemName}

Постои \textbf{поврзан ненасочен едноставен кактус-граф}\textsuperscript{1} со $N \leq 1000$ јазли и $M$ ребра. Неговите јазли имаат бои (означени со ненегативни цели броеви од $0$ до $1499$). Првично, сите јазли имаат боја $0$. \textbf{Детерминистички робот без меморија}\textsuperscript{2} го истражува графот движејќи се од јазол до јазол. Мора да ги посети сите јазли барем еднаш, а потоа да прекине.

Роботот започнува од некој јазол, кој може да биде кој било од јазлите во графот. На секој чекор, ја гледа бојата на својот сегашен јазол и боите на сите соседни јазли \textbf{во некој редослед фиксиран за сегашниот јазол} (т.е. повторното посетување на јазолот ќе му даде на роботот иста низа на соседни јазли, дури и ако нивните бои се различни од претходно). Роботот прави едно од следниве две дејства:
\begin{enumerate}
\item Одлучува да прекине.
\item Избира нова (или можеби иста) боја за сегашниот јазол и на кој соседен јазол да се премести. Соседниот јазол се идентификува со индекс од $0$ до $D-1$, каде што $D$ е бројот на соседни јазли.
\end{enumerate}

Во вториот случај, сегашниот јазол се пребојува (или можеби останува иста боја) и роботот се преместува на избраниот соседен јазол. Ова се повторува сè додека роботот не прекине или додека не ја достигне границата на итерација. Роботот победува ако ги посети сите јазли, а потоа заврши во рамките на ограничувањето на итерацијата од $L = 3000$ чекори (во спротивно губи).

Треба да дизајнирате стратегија за роботот што може да го реши проблемот за кој било кактус-граф. Дополнително, треба да го минимизирате бројот на различни бои што ги користи вашето решение. Бојата 0 секогаш се смета за употребена.

\textsuperscript{1}\textit{Поврзан ненасочен едноставен кактус-граф е поврзан ненасочен едноставен граф (секој јазол е достапен од секој друг јазол; ребрата се двонасочни; нема самојамки или повеќекратни ребра) во кој секое ребро припаѓа на најмногу еден едноставен циклус (едноставен циклус е циклус што го содржи секој јазол најмногу еднаш). Сликата подолу е пример.}

\begin{adjustbox}{center}
\includesvg[inkscapelatex=false, width=.5\linewidth]{graph}
\end{adjustbox}

\textsuperscript{2}\textit{Роботот е детерминистички и без меморија, ако неговото дејство зависи само од неговите сегашни влезни податоци (т.е. не складира податоци од чекор до чекор) и секогаш го избира истото дејство кога му се дадени истите влезни податоци.}

\subsection{Детали за имплементацијата}
Стратегијата на роботот треба да се имплементира како следната функција:
\begin{lstlisting}
std::pair<int, int> navigate(int currColor, std::vector<int> adjColors)
\end{lstlisting}

Како параметри ги прима бојата на сегашниот јазол и боите на сите соседни јазли (по редослед). Мора да врати пар чиј прв елемент е новата боја за сегашниот јазол, а како втор елемент е индексот на соседниот јазол кон кој треба да се движи роботот. Ако наместо тоа, роботот треба да заврши, функцијата треба да го врати парот $(-1, -1)$.

Оваа функција ќе се повикува постојано за да се изберат дејствата на роботот. Бидејќи е детерминистичка, ако \texttt{navigate} веќе е повикана со некои параметри, никогаш повеќе нема да се повика со истите параметри; наместо тоа, нејзината претходна вратена вредност ќе биде повторно употребена. Дополнително, секој тест може да содржи $T \leq 5$ подтестa (различни графови и/или почетни позиции) и тие може да се извршуваат истовремено (т.е. вашата програма може да добива наизменични повици за различни подтести). Конечно, повиците до \texttt{navigate} може да се случат во \textbf{одделни извршувања} на вашата програма (но понекогаш може да се случат и во истото извршување). Вкупниот број на извршувања на вашата програма е $P = 100$. Поради сето ова, вашата програма не треба да се обидува да пренесува информации помеѓу различни повици.

\subsection{Ограничувања}

\begin{itemize}
\item $3 \leq N \leq 1000$
\item $0 \leq \mathrm{Color} < 1500$
\item $L = 3000$
\item $T \leq 5$
\item $P = 100$
\end{itemize}

\subsection{Scoring}

Делот $S$ од поените за подзадача што ќе го добиете зависи од $C$ -- максималниот број на различни бои што ги користи вашето решение (вклучувајќи ја и бојата $0$) на кој било тест во таа подзадача или која било друга задолжителна подзадача:

\begin{itemize}
\item Ако вашето решение не успее на кој било подтест, тогаш $S = 0$.
\item Ако $C е ≤ 4$, тогаш $S = 1.0$.
\item Ако $4 < C е ≤ 8$, тогаш $S = 1.0 - 0.6 \frac{C - 4}{4}$.
\item Ако $8 < C е ≤ 21$, тогаш $S = 0.4 \frac{8}{C}$.
\item Ако $C > 21$, тогаш $S = 0.15$.
\end{itemize}

\subsection{Подзадачи}

\begin{table}[H]
\begin{tblr}{|Q[c,m]|Q[c,m]|Q[c,m]|Q[c,m]|X[c,m]|}
\hline
\textbf{Подзадача} & \textbf{Поени} & \textbf{\makecell{Задолжителни\\подзадачи}} & $N$ & \textbf{Дополнителни ограничувања}\\
\hline
$0$ & $0$ & $-$ & $\leq 300$ & Примерот. \\
\hline
$1$ & $6$ & $-$ & $\leq 300$ & Графот е циклус.\textsuperscript{1} \\
\hline
$2$ & $7$ & $-$ & $\leq 300$ & Графот е ѕвезда.\textsuperscript{2} \\
\hline
$3$ & $9$ & $-$ & $\leq 300$ & Графот е патека.\textsuperscript{3} \\
\hline
$4$ & $16$ & $2 - 3$ & $\leq 300$ & Графот е дрво.\textsuperscript{4} \\
\hline
$5$ & $27$ & $-$ & $\leq 300$ & Сите јазли имаат најмногу $3$ соседни јазли, а јазолот од кој роботот започнува има $1$ соседен јазол. \\
\hline
$6$ & $28$ & $0 - 5$ & $\leq 300$ & $-$ \\
\hline
$7$ & $7$ & $0 - 6$ & $-$ & $-$ \\
\hline
\end{tblr}
\caption*{\textsuperscript{1}Графот на циклусот има рабови: $\left(i, (i+1) \bmod N\right)$ за $0 \leq i < N$. \\
\textsuperscript{2}Графот на ѕвездите има рабови: $\left(0, i\right)$ за $1 \leq i < N$. \\
\textsuperscript{3}Графот на патеките има рабови: $\left(i, i+1\right)$ за $0 \leq i < N - 1$. \\
\textsuperscript{4}Дрвото е граф без циклуси.}
\end{table}
\FloatBarrier

\subsection{Пример}

Разгледајте го примерокот на графикон од сликата во исказот, кој има $N=7$, $M=8$ и рабови $(0, 1)$, $(1, 2)$, $(2, 0)$, $(2, 3)$, $(3, 4)$, $(4, 2)$, $(3, 5)$ и $(2, 6)$. Дополнително, бидејќи редоследот на елементите во листите на соседство на јазлите е релевантен, ги даваме во оваа табела:

\begin{table}[H]
\begin{tblr}{|Q[c,m]|Q[c,m]|}
\hline
\textbf{\makecell{Node}} & \textbf{Соседни јазли} \\
\hline
$0$ & $2, 1$ \\
\hline
$1$ & $2, 0$ \\
\hline
$2$ & $0, 3, 4, 6, 1$ \\
\hline
$3$ & $4, 5, 2$ \\
\hline
$4$ & $2, 3$ \\
\hline
$5$ & $3$ \\
\hline
$6$ & $2$ \\
\hline
\end{tblr}
\end{table}

Да претпоставиме дека роботот започнува на јазолот $5$. Потоа, следново е една можна (неуспешна) низа на интеракции:

\begin{table}[H]
\begin{tblr}{|Q[c,m]|Q[c,m]|Q[c,m]|X[c,m]|Q[c,m]|}
\hline
\textbf{\makecell{\#}} & \textbf{\makecell{Colors}} & \textbf{\makecell{Node}} & \textbf{\makecell{Повик до \texttt{navigate}}} & \textbf{Вратена вредност} \\
\hline
$1$ & $0, 0, 0, 0, 0, 0, 0$ & $5$ & \texttt{navigate(0, \{0\})} & \texttt{\{1, 0\}} \\
\hline
$2$ & $0, 0, 0, 0, 0, 1, 0$ & $3$ & \texttt{navigate(0, \{0, 1, 0\})} & \texttt{\{4, 2\}} \\
\hline
$3$ & $0, 0, 0, 4, 0, 1, 0$ & $2$ & \texttt{navigate(0, \{0, 4, 0, 0, 0\})} & \texttt{\{0, 3\}} \\
\hline
$4$ & $0, 0, 0, 4, 0, 1, 0$ & $6$ & \textsuperscript{1}\texttt{navigate(0, \{0\})} & \texttt{\{1, 0\}} \\
\hline
$5$ & $0, 0, 0, 4, 0, 1, 1$ & $2$ & \texttt{navigate(0, \{0, 4, 0, 1, 0\})} & \texttt{\{8, 0\}} \\
\hline
$6$ & $0, 0, 8, 4, 0, 1, 1$ & $0$ & \texttt{navigate(0, \{8, 0\})} & \texttt{\{3, 0\}} \\
\hline
$7$ & $3, 0, 8, 4, 0, 1, 1$ & $2$ & \texttt{navigate(8, \{3, 4, 0, 1, 0\})} & \texttt{\{2, 2\}} \\
\hline
$8$ & $3, 0, 2, 4, 0, 1, 1$ & $4$ & \texttt{navigate(0, \{2, 4\})} & \texttt{\{1, 1\}} \\
\hline
$9$ & $3, 0, 2, 4, 1, 1, 1$ & $3$ & \texttt{navigate(4, \{1, 1, 2\})} & \texttt{\{-1, -1\}} \\
\hline
\end{tblr}
\end{table}

Тука роботот користел вкупно $6$ различни бои: $0$, $1$, $2$, $3$, $4$ и $8$ (забележете дека $0$ би се сметале за користени дури и ако роботот никогаш не ја вратил бојата $0$, бидејќи сите јазли започнуваат во бојата $0$). Роботот работел $9$ итерации пред да заврши. Сепак, не успеал бидејќи завршил без да го посети јазолот $1$.

\textsuperscript{1}Забележете дека повикот до \texttt{navigate} при итерација $4$ всушност нема да се случи. Ова е затоа што е еквивалентен на повикот при итерација $1$, па затоа оценувачот едноставно ќе ја употреби повратната вредност на вашата функција од тој повик. Сепак, ова сè уште се смета за итерација на роботот.

\subsection{Sample grader}

Пример оценувачот не извршува повеќекратни извршувања на вашата програма, па затоа сите повици до \texttt{navigate} ќе бидат во истото извршување на вашата програма.

Влезниот формат е следниот: Прво се чита $T$ (бројот на подтести). Потоа за секој подтест:
\begin{itemize}
\item line $1$: два цели броја -- $N$ и $M$;
\item line $2+i$ (за $0 \leq i < M$): два цели броја -- $A_i$ и $B_i$, кои се двата јазли што ги поврзува работ $i$ ($0 \leq A_i, B_i < N$).
\end{itemize}

Потоа, пример-оценувачот ќе го испечати бројот на различни бои што ги користело вашето решение и бројот на итерации што му биле потребни пред да заврши. Алтернативно, ќе испечати порака за грешка, ако вашето решение не успее.

Стандардно, пример-оценувачот печати детални информации за тоа што роботот гледа и прави при секоја итерација. Можете да го оневозможите ова, со промена на вредноста на \texttt{DEBUG} од \texttt{true} во \texttt{false}.

\end{document}